\section{问题重述}

% 写作提示：用自己的话交代真实背景、已知条件与数据、待求目标、约束和输出；对含糊概念作必要澄清。
% 不要整段复制题面，也不要在问题重述中提前给出模型公式和求解结果。下方内容仅为结构示例，须按赛题改写。

\subsection{问题背景}
数学建模是解决实际问题的重要工具。本文研究的问题涉及多维度数据分析与优化决策。

\subsection{问题内容}
\textbf{问题一：}对给定数据集进行分析和预测，要求模型具有泛化能力。

\textbf{问题二：}考虑多因素影响，改进和优化模型，提高预测精度。

\textbf{问题三：}结合实际约束条件，设计最优策略方案。

% 技术路线图、研究框架图和算法流程图一律用 paper-diagram Skill 绘制——
% 优先套用其内置模板，产出可编辑 .drawio + PNG，
% 并用 check_layout.py / export_figure.py 校验后插入正文。
